Continuous, wearable seizure monitoring is fundamentally constrained by power, not by detection accuracy in the abstract. This dissertation reports an end-to-end pipeline -- a quantised (w4a4) spatio-temporal convolutional neural network, converted to a spiking neural network and deployed on real BrainChip AKD1000 v1 silicon -- and evaluates it against the CHB-MIT paediatric scalp EEG dataset \citep{goldberger2000} with an explicit emphasis on clinical and chronological honesty throughout.

A full 15-fold leave-one-patient-out (LOPO) cross-validation gives an all-folds mean event sensitivity of 0.586 ($\pm$0.445), a figure shown to be inflated by a 47\% collapse-artefact failure rate; the corrected, collapse-diagnostic-passing figure is 0.470 ($\pm$0.437), against Ali et al.'s unconstrained-compute comparator of 72--75\%. The shape of this result -- markedly bimodal rather than smoothly distributed -- is presented as evidence for genuine inter-patient physiological heterogeneity rather than a simple data-volume shortfall. The dissertation's primary personalisation evidence isolates a single variable -- whether a patient's own data is present in an otherwise near-constant-sized training pool -- across seven patients, finding a real or artefact-clarifying benefit in five of seven cases. A smaller, secondary treatment of per-patient fine-tuning contributes a pool-size finding (collapse-pass rate improving from 37.5\% to 87.5\% with a larger pool) and a conformal-prediction-based correction to an optimistic threshold-selected false-positive-rate figure, presented as a worked example of the project's methodological discipline.

On real AKD1000 v1 hardware, measured energy cost is approximately 0.906\textmu J per inference (chip-isolated, datasheet-derived estimate: approximately 1mW active power), alongside a separately measured, and separately scoped, system-level incremental power draw of 5.63W -- the two figures are reported side by side rather than one validating the other, since they describe genuinely different quantities.

Throughout, this dissertation adopts a triage-device framing, analogous to D-dimer testing in the diagnosis of pulmonary embolism: a low-power, sensitivity-prioritised first-tier screen intended to flag events for closer attention, not a stand-alone replacement for unconstrained-compute clinical systems. Presented alongside its headline results are the project's own caught and corrected methodological errors, treated as evidence of rigour rather than omitted. This work is offered as a first, deliberately modest step toward a genuinely deployable, low-power seizure-screening capability -- one intended, in time and with further validation, to be of real benefit to people living with epilepsy, rather than as a final claim to clinical readiness.
